% Generated from validated Article Draft Result. Do not hand-edit; revise the structured draft instead.
\begingroup\raggedright
\subsection{推論基盤そのものが設計対象になる}
\par\endgroup

3月のsystems系Evidenceでは、KV cache、routing、MoE communication、RPC serving、multimodal serving、memory processing、tool-aware servingが同時に現れる。これはmodelを速くする単一テクニックの集合というより、modelをproduction workloadとしてどう動かすかという層が細分化した姿である。\autocite{src-361b6111931d,src-3802d01433c1,src-446e2c561111,src-795aff00bd48,src-7e670bf82538,src-b08e40540c2a,src-b6e517afaa35,src-d6cda62f02bb,src-fc697f165940}


\begingroup\raggedright
\subsection{TensorRT-LLMとvLLM: stable releaseとpre-releaseを分けて読む}
\par\endgroup

TensorRT-LLMでは3月12日のstable v1.2.0と、月内を通じたv1.3 release-candidate系列を区別する必要がある。v1.3 RCのrelease notesにはdynamic KV-cache quota、reusable KV blocks、prefix-affinity routing、quantizationやMoE servingなど幅広い変更が並ぶが、stable v1.3の成立を意味しない。vLLM v0.18.0ではgRPC servingが追加され、HTTP/REST以外のRPC surfaceもproject側のreleaseとして明示された。\autocite{src-7e670bf82538,src-b08e40540c2a,src-d6cda62f02bb,src-fc697f165940}


\begingroup\raggedright
\subsection{通信・routing・memory・tool executionへ分解されるserving}
\par\endgroup

NCCL EPはauthorsがNCCL Device API上のMoE expert-parallel communication libraryとして提案した。PASTEは繰り返し現れるagent patternから将来のtool invocationを予測し、LLM generationと並行してspeculativeにtoolを実行する構成を提案する。CornserveはAny-to-Any multimodal modelのdistributed servingを、Workload-Router-Pool vision paperはrouting/pool設計を、memory processing論文はPrepare Memory、Compute Relevancy、Retrieval、Apply to Inferenceという4段階の整理を提示した。これらは異なる層だが、推論を単一GPU kernelだけでは捉えにくくなったことを共通して示す。\autocite{src-361b6111931d,src-3802d01433c1,src-446e2c561111,src-795aff00bd48,src-b6e517afaa35}


3月のsystems clusterが厚いのは、agentやmultimodalが新しいworkloadとして現れる一方、その実行コストや通信、memory、routingを受け止める基盤も同時に更新されたためだと解釈できる。PASTEはその接点を明示するが、systems全体をagentの従属テーマへ縮めるべきではない。\autocite{src-361b6111931d,src-446e2c561111,src-d6cda62f02bb,src-fc697f165940}


